\noindent
\begin{minipage}[t]{0.23\textwidth}
\raggedright
\vspace*{530pt}

\footnotesize\sffamily\color{stewartdark}
Ví dụ 1 minh họa rằng ngay cả khi có thể giải một phương trình để tìm $y$ một cách tường minh theo $x$, thì việc sử dụng đạo hàm ẩn vẫn có thể thuận tiện hơn.
\end{minipage}%
\hfill
\begin{minipage}[t]{0.73\textwidth}
\vspace*{0pt}

\noindent\textbf{\textsf{\color{stewartcyan}VÍ DỤ 1}}\hspace{1em}Nếu $x^2 + y^2 = 25$, hãy tìm $\frac{dy}{dx}$. Sau đó, hãy viết phương trình tiếp tuyến của đường tròn $x^2 + y^2 = 25$ tại điểm $(3, 4)$.

\medskip
\noindent\textbf{\textsf{\color{stewartcyan}LỜI GIẢI 1}}\hspace{1em}Lấy đạo hàm cả hai vế của phương trình $x^2 + y^2 = 25$:
\begin{gather*}
\frac{d}{dx}(x^2 + y^2) = \frac{d}{dx}(25) \\[4pt]
\frac{d}{dx}(x^2) + \frac{d}{dx}(y^2) = 0
\end{gather*}
Ghi nhớ rằng $y$ là một hàm số của $x$ và áp dụng Quy tắc chuỗi, ta có
\[
\frac{d}{dx}(y^2) = \frac{d}{dy}(y^2)\,\frac{dy}{dx} = 2y\,\frac{dy}{dx}
\]

\noindent\makebox[0pt][l]{Do đó}\hfill $2x + 2y\,\dfrac{dy}{dx} = 0$\hfill\null
\vspace{2pt}

Bây giờ ta giải phương trình này để tìm $dy/dx$:
\[
\frac{dy}{dx} = -\frac{x}{y}
\]
Tại điểm $(3, 4)$, ta có $x = 3$ và $y = 4$, vì thế
\[
\frac{dy}{dx} = -\frac{3}{4}
\]
Do đó, phương trình tiếp tuyến của đường tròn tại $(3, 4)$ là
\[
y - 4 = -\frac{3}{4}(x - 3) \qquad\text{hay}\qquad 3x + 4y = 25
\]

\medskip
\noindent\textbf{\textsf{\color{stewartcyan}LỜI GIẢI 2}}\hspace{1em}Giải phương trình $x^2 + y^2 = 25$ theo $y$, ta nhận được $y = \pm\sqrt{25 - x^2}$. Điểm $(3, 4)$ nằm trên nửa đường tròn phía trên $y = \sqrt{25 - x^2}$ và do đó ta xét hàm số $f(x) = \sqrt{25 - x^2}$. Lấy đạo hàm của $f$ nhờ sử dụng Quy tắc chuỗi, ta có
\begin{align*}
f'(x) &= \frac{1}{2}(25 - x^2)^{-1/2}\,\frac{d}{dx}(25 - x^2) \\[6pt]
&= \frac{1}{2}(25 - x^2)^{-1/2}(-2x) = -\frac{x}{\sqrt{25 - x^2}}
\end{align*}
Tại điểm $(3, 4)$, ta có
\[
f'(3) = -\frac{3}{\sqrt{25 - 3^2}} = -\frac{3}{4}
\]
và tương tự như trong Lời giải 1, phương trình của tiếp tuyến là $3x + 4y = 25$. \hfill{\color{stewartcyan}$\blacksquare$}

\bigskip
\noindent\textbf{\textsf{\color{stewartgreen}LƯU Ý 1}}\hspace{1em}Biểu thức $dy/dx = -x/y$ trong Lời giải 1 cho đạo hàm theo cả $x$ và $y$. Biểu thức này luôn đúng bất kể hàm số $y$ nào được xác định bởi phương trình đã cho. Chẳng hạn, với $y = f(x) = \sqrt{25 - x^2}$, ta có
\[
\frac{dy}{dx} = -\frac{x}{y} = -\frac{x}{\sqrt{25 - x^2}}
\]

\end{minipage}
